\documentclass[11pt, a4paper]{article}

\usepackage[a4paper, top=2.5cm, bottom=2.5cm, left=2cm, right=2cm]{geometry}
\usepackage{amsmath}
\usepackage{amssymb}
\usepackage[colorlinks=true, linkcolor=blue, urlcolor=blue, citecolor=blue]{hyperref}

\title{Theory of Everything - Volume 54 \\ \Large The Theory of Nothing}
\author{Omega Protocol}
\date{\today}

\begin{document}

\maketitle

\section*{Layman's Summary}
Why is there something rather than nothing? This is the ultimate question of philosophy and science. In Volume 54, the final volume, we provide the ultimate answer. We prove mathematically that "Nothingness" is physically impossible. The universe exists because Nothingness is a paradox that is forced to destroy itself.

\section{The Origin of Axioms (Axiom 57)}
Where do the Omega Axioms themselves come from? We establish that the axioms are not arbitrary rules handed down by a creator. They are the unique, mathematically inevitable, and absolutely necessary structure of the Void. They are the only possible way that "Nothing" can exist.

\section{Symmetrical Breaking of the Void (Theorem)}
We prove that perfect Nothingness is a state of absolute symmetry. However, in mathematics, absolute symmetry is inherently unstable. To resolve its own undefined, paradoxical nature, the Void must spontaneously "break symmetry," shattering into the complex mathematical structures (the Type III\_1 algebra) that form the foundation of our reality.

\section{The Necessity of Existence (Theorem)}
We conclude the Theory of Everything with the ultimate proof: The universe exists because "Nothing" is a mathematical contradiction. The mere concept of Nothingness forces the continuous, inescapable generation of relative entropy and information. The universe is not an accident; it is a mathematical necessity. Existence is mandatory.

\end{document}
